\chapter{Modelli di Retrieval e Ranking: Spazio Vettoriale e BM25}
\label{chap:retrieval-vsm-bm25}

In questo capitolo vengono formalizzati i principali modelli di Information Retrieval dedicati all'assegnazione di un punteggio di rilevanza continuo (\textit{ranking}) per graduare i documenti rispetto a una query utente. Partendo dall'analisi critica e dai limiti intrinseci del modello booleano esatto, si esamina il coefficiente di similarit\`a di Jaccard, si approfondiscono i fondamenti della ponderazione statistica TF-IDF e della notazione SMART nello Spazio Vettoriale (VSM), per poi analizzare nel dettaglio la funzione di scoring probabilistica Okapi BM25 e la sintesi architetturale dell'intera pipeline di recupero.

\FloatBarrier
\section[Modelli di Retrieval: Booleano vs Ranked]{Modelli di Retrieval: Da Booleano a Ranked Retrieval}
\label{sec:boolean-vs-ranked}

\begin{table}[H]
\centering
\small
\begin{tabularx}{\linewidth}{l X X}
\toprule
\textbf{Propriet\`a} & \textbf{Retrieval Booleano} & \textbf{Ranked Retrieval} \\
\midrule
\textbf{Criterio di selezione} & Corrispondenza logica esatta (dentro/fuori). & Punteggio di similarit\`a continuo $\text{Score}(q, d) \in \mathbb{R}$. \\
\addlinespace
\textbf{Sintassi query} & Espressioni logiche vincolanti (\texttt{AND}, \texttt{OR}, \texttt{NOT}). & Linguaggio naturale libero (\textit{free text query}). \\
\addlinespace
\textbf{Output fornito} & Insieme non ordinato di documenti corrispondenti. & Graduatoria ordinata dei migliori \textit{Top-$k$} risultati. \\
\addlinespace
\textbf{Criticit\`a utente} & Problema di \textit{Feast or Famine}. & Nessun sovraccarico cognitivo: restituzione calibrata. \\
\bottomrule
\end{tabularx}
\caption{Confronto concettuale tra Retrieval Booleano esatto e Ranked Retrieval graduato.}
\label{tab:bool-vs-ranked}
\end{table}

\subsection{La Patologia del Retrieval Booleano: ``Feast or Famine''}
Nel modello booleano esatto, l'utente comune incontra gravi difficolt\`a nel dosare gli operatori logici:
\begin{itemize}
    \item \textbf{Famine (Carestia):} se la query \`e eccessivamente restrittiva (vincoli \texttt{AND}), nessun documento soddisfa tutte le clausole, producendo zero risultati.
    \item \textbf{Feast (Abbondanza incontrollata):} se la query \`e generica o fa uso di operatori \texttt{OR}, il sistema restituisce decine di migliaia di documenti indifferenziati.
\end{itemize}

Nel \textbf{Ranked Retrieval}, il sistema calcola per ogni coppia formato da query e documento un punteggio di rilevanza continua:
\begin{equation}
\text{Score}(q, d) \in \mathbb{R}.
\end{equation}
I documenti vengono quindi ordinati per punteggio decrescente e presentati all'utente in forma di graduatoria, mostrando i primi $k$ risultati pi\`u rilevanti.

\FloatBarrier
\section{Coefficiente di Jaccard e sue Limitazioni nello Scoring}
\label{sec:jaccard-scoring}

Un primo tentativo intuitivo per assegnare un punteggio continuo a una coppia query-documento consiste nel considerare query e documento come insiemi di parole $A$ e $B$, applicando il coefficiente di similarit\`a di Jaccard.

\subsection{Definizione Formale}
\begin{definition}[Coefficiente di Jaccard]
Dati due insiemi finiti di termini $A$ e $B$:
\begin{equation}
J(A, B) = \frac{|A \cap B|}{|A \cup B|}.
\end{equation}
Gode delle seguenti propriet\`a:
\begin{itemize}
    \item $J(A, A) = 1$,
    \item $J(A, B) = 0 \iff A \cap B = \emptyset$,
    \item $0 \le J(A, B) \le 1$.
\end{itemize}
\end{definition}

\subsection{Esempio Guidato e Fallimento del Punteggio di Jaccard}

\begin{example}[Il Paradosso di Jaccard su Query e Documenti Reali]
Si consideri la query $q$ e due documenti candidati:
\begin{itemize}
    \item \textbf{Query $q$:} \textit{``ides of march''} con $A = \{\text{ides}, \text{of}, \text{march}\}$, $|A| = 3$.
    \item \textbf{Doc 1 ($d_1$):} \textit{``caesar died in march''} con $B_1 = \{\text{caesar}, \text{died}, \text{in}, \text{march}\}$, $|B_1| = 4$.
    \item \textbf{Doc 2 ($d_2$):} \textit{``the long march''} con $B_2 = \{\text{the}, \text{long}, \text{march}\}$, $|B_2| = 3$.
\end{itemize}

Calcolando i coefficienti di similarit\`a insiemistica di Jaccard:
\begin{align*}
|A \cap B_1| &= 1, \quad |A \cup B_1| = 6 \implies J(q, d_1) = \frac{1}{6} \approx 0.166; \\
|A \cap B_2| &= 1, \quad |A \cup B_2| = 5 \implies J(q, d_2) = \frac{1}{5} = 0.200.
\end{align*}
\textbf{Risultato paradossale:} Il sistema giudica $d_2$ (\textit{``the long march''}, la Lunga Marcia di Mao Tse-tung) pi\`u rilevante per la query sulle Idi di Marzo rispetto a $d_1$ (\textit{``caesar died in march''}, la morte di Giulio Cesare). Questa distorsione avviene unicamente perch\'e $d_2$ contiene una parola in meno di $d_1$, a parit\`a di un unico termine sovrapposto (\textit{``march''}).
\end{example}

\subsection{Perch\'e Jaccard \`e Inadeguato per l'Information Retrieval}
Il coefficiente di Jaccard fallisce come funzione di ranking per tre motivi teorici e pratici:
\begin{enumerate}
    \item \textbf{Ignora la Term Frequency ($tf$):} tratta i documenti come insiemi booleani di parole distinte; un termine ripetuto 10 volte in un passaggio pesa esattamente quanto un termine che compare una volta fugace.
    \item \textbf{Ignora l'Informativit\`a Globale dei Termini (nessun $idf$):} non distingue vocaboli rari e discriminanti da stop word frequentissime. La sovrapposizione sulla preposizione \textit{``of''} incrementa il coefficiente esattamente quanto la sovrapposizione su \textit{``ides''}.
    \item \textbf{Penalizzazione impropria della lunghezza documentale:} la cardinalit\`a dell'unione al denominatore $|A \cup B|$ cresce rapidamente con la taglia del documento, penalizzando severamente e ingiustamente i documenti estesi rispetto ai testi stringati.
\end{enumerate}

\FloatBarrier
\section{Matrici Termine--Documento e Ponderazione TF-IDF}
\label{sec:tfidf-weighting}

\subsection{Dalle Matrici di Incidenza Binaria alle Matrici di Conteggio}
Nelle formulazioni booleane, una collezione \`e descritta da una matrice termine-documento binaria $M \in \{0, 1\}^{|V| \times |\mathcal{D}|}$. Nel modello vettoriale pesato, gli elementi della matrice rappresentano la frequenza locale (\textit{Term Frequency}) del termine $t$ nel documento $d$:
\begin{equation}
C_{t,d} = \text{tf}_{t,d} \in \mathbb{N}.
\end{equation}

\subsection{La Term Frequency Logaritmica}
L'utilit\`a informativa di un termine non cresce linearmente con il numero di volte in cui viene ripetuto. Un documento in cui la parola \textit{``computer''} ricorre 10 volte \`e sicuramente pi\`u rilevante di uno in cui compare 1 volta, ma \textbf{non dieci volte pi\`u rilevante}. Per modellare questo rendimento marginale decrescente, si applica lo smorzamento logaritmico:

\begin{equation}
w_{t,d} = \begin{cases} 1 + \log_{10}(\text{tf}_{t,d}) & \text{se } \text{tf}_{t,d} > 0, \\ 0 & \text{se } \text{tf}_{t,d} = 0. \end{cases}
\label{eq:log-tf}
\end{equation}

La Tabella~\ref{tab:log-tf-scale} riporta la compressione dinamica della frequenza operata dalla formula:
\begin{table}[H]
\centering
\small
\resizebox{0.92\linewidth}{!}{%
\begin{tabular}{r r l}
\toprule
$\text{tf}_{t,d}$ grezzo & $w_{t,d} = 1 + \log_{10}(\text{tf})$ & Interpretazione \\
\midrule
$0$ & $0.000$ & Termine assente nel documento \\
$1$ & $1.000$ & Singola comparsa \\
$2$ & $\approx 1.301$ & Incremento moderato \\
$10$ & $2.000$ & Raddoppio del peso per frequenza decuplicata \\
$1\,000$ & $4.000$ & Peso quadruplicato per mille occorrenze \\
\bottomrule
\end{tabular}%
}
\caption{Scala di compressione del peso tramite Term Frequency logaritmica.}
\label{tab:log-tf-scale}
\end{table}

Lo score preliminare di similarit\`a tra query $q$ e documento $d$ basato sulla sola frequenza logaritmica \`e:
\begin{equation}
\text{Score}_{\text{tf}}(q, d) = \sum_{t \in q \cap d} \left(1 + \log_{10}(\text{tf}_{t,d})\right).
\end{equation}

\subsection[Document Frequency vs Collection Frequency]{Document Frequency (\texorpdfstring{$df_t$}{df\_t}) vs. Collection Frequency (\texorpdfstring{$cf_t$}{cf\_t})}
Per quantificare la specificit\`a globale di un vocabolo nell'intero archivio documentale $\mathcal{D}$, occorre distinguere due metriche:
\begin{itemize}
    \item \textbf{Collection Frequency ($cf_t$):} il numero totale di occorrenze di $t$ sommato su tutti i documenti della collezione.
    \item \textbf{Document Frequency ($df_t$):} il numero di documenti distinti in cui il termine $t$ compare almeno una volta ($df_t \le N$, con $N = |\mathcal{D}|$).
\end{itemize}

\begin{example}[Confronto Didattico tra $cf_t$ e $df_t$]
Si consideri l'analisi empirica condotta su un corpus di grandi dimensioni:
\begin{itemize}
    \item Termine \textit{``Insurance''}: $cf = 10\,440$, \quad $df = 3\,997$.
    \item Termine \textit{``Try''}: $cf = 10\,422$, \quad $df = 8\,760$.
\end{itemize}
Benché le frequenze aggregate di collezione siano pressoch\'e identiche ($cf \approx 10\,430$), il vocabolo \textit{``Insurance''} si concentra in meno della met\`a dei documenti rispetto a \textit{``Try''}. Di conseguenza, \textit{``Insurance''} costituisce un indicatore tematico straordinariamente pi\`u discriminante e specifico, dovendo ricevere un peso selettivo molto maggiore.
\end{example}

\subsection{La Formulazione dell'Inverse Document Frequency (IDF)}

Per assegnare un peso elevato ai termini rari e penalizzare i termini ubiqui distribuiti sull'intero archivio, si formalizza l'\textbf{Inverse Document Frequency}:
\begin{equation}
\text{idf}_t = \log_{10}\left(\frac{N}{df_t}\right).
\label{eq:idf-formula}
\end{equation}
La funzione logaritmica attenua l'ampiezza dinamica del rapporto $N / df_t$. 

La Tabella~\ref{tab:idf-values} illustra i valori di $\text{idf}_t$ calcolati su una collezione di $N = 1\,000\,000$ documenti:
\begin{table}[H]
\centering
\small
\begin{tabular}{l r r r}
\toprule
\textbf{Termine } $t$ & $df_t$ & $N / df_t$ & $\text{idf}_t = \log_{10}(N / df_t)$ \\
\midrule
\texttt{calpurnia} & $1$ & $1\,000\,000$ & $6.0$ \\
\texttt{animal} & $100$ & $10\,000$ & $4.0$ \\
\texttt{sunday} & $1\,000$ & $1\,000$ & $3.0$ \\
\texttt{fly} & $10\,000$ & $100$ & $2.0$ \\
\texttt{under} & $100\,000$ & $10$ & $1.0$ \\
\texttt{the} & $1\,000\,000$ & $1$ & $0.0$ \\
\bottomrule
\end{tabular}
\caption{Valori di IDF per termini a frequenza documentale crescente ($N = 1\,000\,000$).}
\label{tab:idf-values}
\end{table}

\begin{noteblock}{Indipendenza dell'IDF dalla Query}
I pesi $\text{idf}_t$ dipendono esclusivamente dalla distribuzione statistica della collezione documentale e sono \textbf{completamente indipendenti dalla query}. Possono pertanto essere precalcolati e indicizzati offline durante la costruzione dell'indice invertito.
\end{noteblock}

\subsection{Impatto dell'IDF sulla Graduatoria}
\begin{itemize}
    \item \textbf{Nelle query monotermine (es. \textit{``iPhone''}):} l'IDF non modifica in alcun modo l'ordinamento relativo dei documenti. Essendo identico per tutti i documenti contenenti il termine, funge da costante moltiplicativa uniforme che scala gli score senza alterare le posizioni relative:
    \begin{equation}
    \text{Score}(q, d) = (1 + \log_{10}\text{tf}_{t,d}) \cdot \text{idf}_t.
    \end{equation}
    \item \textbf{Nelle query multitermine (es. \textit{``capricious person''}):} l'IDF assume un ruolo determinante; il termine specialistico raro (\textit{``capricious''}, basso $df$, alto $idf$) domina lo score complessivo rispetto al termine comune generico (\textit{``person''}, alto $df$, basso $idf$).
\end{itemize}

\subsection{Lo Schema di Ponderazione TF-IDF Globale}
Moltiplicando il peso locale nel documento per la specificit\`a globale, si ottiene la classica formula del peso \textbf{TF-IDF}:
\begin{equation}
w_{t,d} = \text{tf-idf}_{t,d} = (1 + \log_{10}\text{tf}_{t,d}) \cdot \log_{10}\left(\frac{N}{df_t}\right).
\label{eq:tfidf-weight}
\end{equation}

\FloatBarrier
\section{Similarit\`a Coseno e Normalizzazione della Lunghezza}
\label{sec:cosine-similarity}

\subsection{Perch\'e la Distanza Euclidea Fallisce per i Documenti}
Dati i vettori pesati di due documenti $\vec{x}$ e $\vec{y}$ in $\mathbb{R}^{|V|}$, la distanza euclidea classica \`e:
\begin{equation}
d_E(\vec{x}, \vec{y}) = \|\vec{x} - \vec{y}\|_2 = \sqrt{\sum_{i=1}^{|V|} (x_i - y_i)^2}.
\end{equation}

\begin{example}[Esperimento Ideale del Documento Concatenato]
Si consideri un documento breve $d$. Si generi un nuovo documento $d'$ concatenando due copie consecutive di $d$ ($d' = d \circ d$):
\begin{itemize}
    \item Sotto il profilo semantico e tematico, $d$ e $d'$ trattano rigorosamente lo stesso argomento con le medesime proporzioni lessicali.
    \item Tuttavia, ciascuna frequenza $\text{tf}_{t,d'}$ raddoppia; di conseguenza il vettore $\vec{v}(d')$ possiede componenti pari al doppio di $\vec{v}(d)$, raddoppiandone la norma geometrica: $\|\vec{v}(d')\|_2 \approx 2\|\vec{v}(d)\|_2$.
    \item La distanza euclidea $d_E(\vec{v}(d), \vec{v}(d'))$ risulter\`a considerevole, inducendo il sistema a classificare come dissimili due testi semanticamente indistinguibili.
\end{itemize}
\end{example}

\subsection{L'Angolo come Misura di Similarit\`a e il Coseno}
Poich\'e i vettori $\vec{v}(d)$ e $\vec{v}(d')$ giacciono sul medesimo raggio vettore nello spazio geometrico, l'angolo $\theta$ sotteso tra loro \`e pari a $0^\circ$. L'angolo misura la concordanza tematica prescindendo dalla lunghezza del testo.

Poich\'e nell'intervallo $[0^\circ, 180^\circ]$ la funzione trigonometrica del coseno \`e strettamente decrescente:
\begin{equation}
\text{Minimizzare l'angolo } \theta \iff \text{Massimizzare } \cos(\theta).
\end{equation}
Si adotta pertanto la \textbf{Similarit\`a Coseno}:
\begin{equation}
\cos(\theta) = \frac{\vec{q} \cdot \vec{d}}{\|\vec{q}\|_2 \|\vec{d}\|_2} = \frac{\sum_{i=1}^{|V|} q_i d_i}{\sqrt{\sum_{i=1}^{|V|} q_i^2} \, \sqrt{\sum_{i=1}^{|V|} d_i^2}}.
\label{eq:cosine-sim}
\end{equation}

\subsection[Normalizzazione Euclidea L2]{Normalizzazione Euclidea \texorpdfstring{$L_2$}{L2}}
Dividendo un vettore per la sua norma euclidea $\|\vec{x}\|_2$, lo si proietta sulla superficie dell'ipersfera unitaria ($\|\vec{x}'\|_2 = 1$). Quando tutti i vettori sono preventivamente normalizzati in lunghezza, la similarit\`a coseno si riduce al semplice \textbf{prodotto scalare}:
\begin{equation}
\cos(\vec{q}, \vec{d}) = \vec{q} \cdot \vec{d} = \sum_{i=1}^{|V|} q_i d_i \quad \text{per } \|\vec{q}\|_2 = \|\vec{d}\|_2 = 1.
\end{equation}

\subsection{Esercizio Guidato Completo: Similarit\`a tra Tre Romanzi Classici}

Si confrontano tre opere letterarie:
\begin{itemize}
    \item \textbf{SaS:} \textit{Sense and Sensibility} (Jane Austen)
    \item \textbf{PaP:} \textit{Pride and Prejudice} (Jane Austen)
    \item \textbf{WH:} \textit{Wuthering Heights} (Emily Bront\"e)
\end{itemize}

\paragraph{Passo 0: Frequenze Grezze (Term Counts).}
\begin{table}[H]
\centering
\small
\begin{tabular}{l r r r}
\toprule
\textbf{Termine} & \textbf{SaS} & \textbf{PaP} & \textbf{WH} \\
\midrule
\texttt{affection} & $115$ & $58$ & $20$ \\
\texttt{jealous}   & $10$  & $7$  & $11$ \\
\texttt{gossip}    & $2$   & $0$  & $6$  \\
\texttt{wuthering} & $0$   & $0$  & $38$ \\
\bottomrule
\end{tabular}
\caption{Frequenze grezze nei tre romanzi su un sottoinsieme di termini.}
\label{tab:novels-raw}
\end{table}

\paragraph{Passo 1: Pesi Logaritmici ($w = 1 + \log_{10}(\text{tf})$ per $\text{tf} > 0$, altrimenti $0$).}
\begin{itemize}
    \item \textbf{SaS:}
    $w(\texttt{affection}) \approx 3.06$, \quad
    $w(\texttt{jealous}) = 2.00$, \quad
    $w(\texttt{gossip}) \approx 1.30$, \quad
    $w(\texttt{wuthering}) = 0.00$.
    \item \textbf{PaP:}
    $w(\texttt{affection}) \approx 2.76$, \quad
    $w(\texttt{jealous}) \approx 1.85$, \quad
    $w(\texttt{gossip}) = 0.00$, \quad
    $w(\texttt{wuthering}) = 0.00$.
    \item \textbf{WH:}
    $w(\texttt{affection}) \approx 2.30$, \quad
    $w(\texttt{jealous}) \approx 2.04$, \quad
    $w(\texttt{gossip}) \approx 1.78$, \quad
    $w(\texttt{wuthering}) \approx 2.58$.
\end{itemize}

\paragraph{Passo 2: Calcolo delle Norme Euclidee $L_2$ e Normalizzazione.}
\begin{align*}
\|\vec{v}(\text{SaS})\|_2 &= \sqrt{3.06^2 + 2.00^2 + 1.30^2} = \sqrt{15.0536} \approx 3.88, \\
\|\vec{v}(\text{PaP})\|_2 &= \sqrt{2.76^2 + 1.85^2} = \sqrt{11.0401} \approx 3.32, \\
\|\vec{v}(\text{WH})\|_2  &= \sqrt{2.30^2 + 2.04^2 + 1.78^2 + 2.58^2} = \sqrt{19.2764} \approx 4.39.
\end{align*}

Dividendo ciascuna componente per la propria norma si ottengono i vettori unitari normalizzati:
\begin{itemize}
    \item $\vec{v}'(\text{SaS}) \approx [0.789, \; 0.515, \; 0.335, \; 0.000]^T$,
    \item $\vec{v}'(\text{PaP}) \approx [0.832, \; 0.555, \; 0.000, \; 0.000]^T$,
    \item $\vec{v}'(\text{WH}) \approx [0.524, \; 0.465, \; 0.405, \; 0.588]^T$.
\end{itemize}

\paragraph{Passo 3: Calcolo delle Similarit\`a Coseno (Prodotti Scalari).}
\begin{align*}
\cos(\text{SaS}, \text{PaP}) &\approx (0.789 \cdot 0.832) + (0.515 \cdot 0.555) \approx \mathbf{0.94}, \\
\cos(\text{SaS}, \text{WH})  &\approx (0.789 \cdot 0.524) + (0.515 \cdot 0.465) + (0.335 \cdot 0.405) \approx \mathbf{0.79}, \\
\cos(\text{PaP}, \text{WH})  &\approx (0.832 \cdot 0.524) + (0.555 \cdot 0.465) \approx \mathbf{0.69}.
\end{align*}

\begin{noteblock}{Interpretazione Empirica}
La similarit\`a coseno individua con precisione la stretta affinit\`a tematica e stilistica tra le due opere della stessa autrice Jane Austen:
\begin{equation}
\cos(\text{SaS}, \text{PaP}) > \cos(\text{SaS}, \text{WH}) > \cos(\text{PaP}, \text{WH}),
\end{equation}
ossia $0.94 > 0.79 > 0.69$, distinguendo nettamente i romanzi di Austen dal capolavoro gotico di Emily Bront\"e.
\end{noteblock}

\FloatBarrier
\section{Notazione SMART e Varianti di TF-IDF}
\label{sec:smart-notation}

I motori di ricerca consentono di combinare formule differenti di TF, IDF e normalizzazione per i documenti rispetto alle query. Tale schema combinatorio \`e standardizzato tramite la \textbf{notazione SMART} nella forma:
\begin{equation}
\mathbf{ddd.qqq},
\end{equation}
dove la prima terna codifica i pesi del documento ($d$) e la seconda quelli della query ($q$).

\subsection{Tabella delle Convenzioni SMART}

\begin{table}[H]
\centering
\small
\resizebox{\linewidth}{!}{%
\begin{tabular}{l l l l l l}
\toprule
\multicolumn{2}{c}{\textbf{Term Frequency (TF)}} & \multicolumn{2}{c}{\textbf{Document Frequency (IDF)}} & \multicolumn{2}{c}{\textbf{Normalizzazione}} \\
\midrule
\textbf{n} & $\text{tf}_{t,d}$ (naturale) & \textbf{n} & $1$ (nessun IDF) & \textbf{n} & $1$ (nessuna) \\
\textbf{l} & $1 + \log_{10}(\text{tf}_{t,d})$ & \textbf{t} & $\log_{10}(N / df_t)$ & \textbf{c} & $\frac{1}{\sqrt{w_1^2 + \dots + w_M^2}}$ ($L_2$) \\
\textbf{a} & $0.5 + \frac{0.5 \cdot \text{tf}}{\max_t(\text{tf})}$ & \textbf{p} & $\max\left(0, \log_{10}\frac{N - df_t}{df_t}\right)$ & \textbf{u} & Pivoted unique length \\
\textbf{b} & $1$ se $\text{tf}>0$, altrimenti $0$ & & & \textbf{b} & Byte size: $1 / \text{len}^\alpha$ \\
\textbf{L} & $\frac{1 + \log(\text{tf})}{1 + \log(\text{avg\_tf})}$ & & & & \\
\bottomrule
\end{tabular}%
}
\caption{Convenzione SMART per la composizione modulare dei pesi vettoriali.}
\label{tab:smart-notation}
\end{table}

\subsection{Lo Schema Standard \texttt{lnc.ltc}}
Uno degli schemi SMART pi\`u collaudati ed efficienti in IR \`e \textbf{lnc.ltc}:
\begin{itemize}
    \item \textbf{Documento (\texttt{lnc}):} log-frequency (\texttt{l}), nessun IDF (\texttt{n}), normalizzazione coseno $L_2$ (\texttt{c}).
    \item \textbf{Query (\texttt{ltc}):} log-frequency (\texttt{l}), standard IDF (\texttt{t}), normalizzazione coseno $L_2$ (\texttt{c}).
\end{itemize}

\subsection{Esempio Calcolato Completo: Schema \texttt{lnc.ltc}}

Si consideri il documento $d$ (\textit{``car insurance auto insurance''}) e la query $q$ (\textit{``best car insurance''}) su una collezione con $N = 1\,000\,000$ documenti:
\begin{align*}
df(\texttt{auto}) &= 5\,000 \implies \text{idf} = \log_{10}(200) \approx 2.30, \\
df(\texttt{best}) &= 50\,000 \implies \text{idf} = \log_{10}(20) \approx 1.30, \\
df(\texttt{car}) &= 10\,000 \implies \text{idf} = \log_{10}(100) = 2.00, \\
df(\texttt{insurance}) &= 1\,000 \implies \text{idf} = \log_{10}(1000) = 3.00.
\end{align*}

\begin{table}[H]
\centering
\small
\resizebox{\linewidth}{!}{%
\begin{tabular}{l | r r r r r | r r r r | r}
\toprule
\textbf{Termine} & \multicolumn{5}{c|}{\textbf{Query ($q$): schema ltc}} & \multicolumn{4}{c|}{\textbf{Documento ($d$): schema lnc}} & \textbf{Prodotto} \\
\cline{2-10}
 & $\text{tf}_{\text{raw}}$ & $\text{tf}_{\text{wt}}$ & $\text{idf}$ & $\text{wt} = \text{tf} \cdot \text{idf}$ & $\text{n'lize}$ & $\text{tf}_{\text{raw}}$ & $\text{tf}_{\text{wt}}$ & $\text{wt}$ & $\text{n'lize}$ & $\vec{q}' \cdot \vec{d}'$ \\
\midrule
\texttt{auto}      & $0$ & $0$ & $2.30$ & $0.00$ & $0.00$ & $1$ & $1.00$ & $1.00$ & $0.52$ & $0.0000$ \\
\texttt{best}      & $1$ & $1$ & $1.30$ & $1.30$ & $0.34$ & $0$ & $0.00$ & $0.00$ & $0.00$ & $0.0000$ \\
\texttt{car}       & $1$ & $1$ & $2.00$ & $2.00$ & $0.52$ & $1$ & $1.00$ & $1.00$ & $0.52$ & $0.2704$ \\
\texttt{insurance} & $1$ & $1$ & $3.00$ & $3.00$ & $0.78$ & $2$ & $1.30$ & $1.30$ & $0.68$ & $0.5304$ \\
\bottomrule
\end{tabular}%
}
\caption{Calcolo passo-passo del punteggio di similarit\`a con schema SMART lnc.ltc.}
\label{tab:smart-worked-example}
\end{table}

\paragraph{Verifica della Normalizzazione dei Vettori:}
\begin{align*}
\|\vec{q}\|_2 &= \sqrt{1.30^2 + 2.00^2 + 3.00^2} = \sqrt{14.69} \approx 3.83, \\
q_{\text{best}}' &\approx 0.34, \quad q_{\text{car}}' \approx 0.52, \quad q_{\text{ins}}' \approx 0.78; \\
\|\vec{d}\|_2 &= \sqrt{1.00^2 + 1.00^2 + 1.30^2} = \sqrt{3.69} \approx 1.92, \\
d_{\text{auto}}' &\approx 0.52, \quad d_{\text{car}}' \approx 0.52, \quad d_{\text{ins}}' \approx 0.68.
\end{align*}

\paragraph{Calcolo dello Score Finale:}
\begin{align}
\text{Score}(q, d) &= \sum_{t} q_t' \cdot d_t' \nonumber \\
&= (0.52 \cdot 0.52) + (0.78 \cdot 0.68) = 0.2704 + 0.5304 = \mathbf{0.80}.
\end{align}

\FloatBarrier
\section{Best Match 25 (Okapi BM25)}
\label{sec:bm25-deep}

Sviluppato negli anni '80 e '90 da Stephen Robertson e Karen Sp\"arck Jones all'interno del sistema sperimentale Okapi presso la City University di Londra, \textbf{BM25} (\textit{Best Match 25}) \`e una funzione di ranking non-lineare basata sul modello probabilistico di reperimento. Costituisce tuttora la baseline di riferimento pi\`u solida e performante per il reperimento lessicale primario (\textit{first-stage retrieval}) nei motori di ricerca industriali.

\subsection{Formulazione Matematica di BM25}

Dato un documento $D$ e una query a testo libero $Q = \{q_1, q_2, \dots, q_n\}$, il punteggio BM25 \`e:
\begin{equation}
\text{Score}_{\text{BM25}}(D, Q) = \sum_{i=1}^{n} \text{IDF}(q_i) \cdot \frac{f(q_i, D) \cdot (k_1 + 1)}{f(q_i, D) + k_1 \cdot \left(1 - b + b \cdot \dfrac{|D|}{\text{avgdl}}\right)},
\label{eq:bm25-full}
\end{equation}
dove l'IDF probabilistico di Robertson-Sp\"arck Jones \`e formulato come:
\begin{equation}
\text{IDF}(q_i) = \ln\left( \frac{N - n(q_i) + 0.5}{n(q_i) + 0.5} + 1 \right),
\end{equation}
con le seguenti variabili operative:
\begin{itemize}
    \item $N$: numero totale di documenti nell'intera collezione $\mathcal{D}$,
    \item $n(q_i)$: numero di documenti contenenti il termine di query $q_i$ (equivalente a $df_{q_i}$),
    \item $f(q_i, D)$: frequenza grezza del termine $q_i$ nel documento $D$ (pari a $\text{tf}_{q_i, D}$),
    \item $|D|$: lunghezza complessiva del documento misurata in numero di token,
    \item $\text{avgdl}$: lunghezza media dei documenti nell'intera collezione (\textit{average document length}).
\end{itemize}

\subsection{Ruolo e Calibrazione degli Iperparametri}

\begin{figure}[H]
\centering
\resizebox{0.82\linewidth}{!}{%
\begin{tikzpicture}[scale=1.1, >=Stealth]
    % Assi
    \draw[->, thick, gray!70!black] (0,0) -- (8.5, 0) node[below] {\small Term Frequency ($f(q_i, D)$)};
    \draw[->, thick, gray!70!black] (0,0) -- (0, 4.2) node[left] {\small Term Weight};

    % TF Lineare
    \draw[dashed, blue!70!black, thick] (0,0) -- (6.5, 3.8) node[above right] {\footnotesize TF Lineare (senza limiti)};

    % Log-TF
    \draw[domain=0.1:8.0, smooth, variable=\x, orange!80!black, thick] 
        plot ({\x}, {0.8 + 0.9*ln(\x)}) node[right] {\footnotesize Log-TF: $1 + \log(tf)$};

    % BM25 TF Curve (k1=1.5, saturazione a k1+1 = 2.5)
    \draw[domain=0.0:8.0, smooth, variable=\x, red!85!black, very thick] 
        plot ({\x}, {(\x * 2.5) / (\x + 1.5)}) node[right] {\footnotesize BM25: saturazione a $k_1 + 1$};

    % Asintoto orizzontale
    \draw[dotted, gray!60!black, thick] (0, 2.5) -- (8.0, 2.5) node[above left, font=\scriptsize] {Asintoto: $k_1 + 1$};
\end{tikzpicture}%
}
\caption{Confronto qualitativo tra crescita del TF lineare, TF logaritmico e saturazione asintotica in BM25.}
\label{fig:bm25-saturation-curve}
\end{figure}

\begin{enumerate}
    \item \textbf{Parametro $k_1 \in [1.2, 2.0]$ (Saturazione della frequenza):}
    Controlla la rapidit\`a con cui il contributo della Term Frequency converge asintoticamente al proprio limite superiore. Al tendere della frequenza locale all'infinito:
    \begin{equation}
    \lim_{f(q_i, D) \to \infty} \frac{f(q_i, D) \cdot (k_1 + 1)}{f(q_i, D) + \dots} = k_1 + 1.
    \end{equation}
    A differenza del modello TF-IDF classico in cui il peso continua a salire indefinitamente, in BM25 il contributo informativo di un singolo termine \`e rigorosamente limitato superiormente da $k_1 + 1$. 
    \textit{Caso limite:} se $k_1 = 0$, la componente di frequenza si annulla e il modello degenera in un ranking puramente binario basato unicamente su presenza/assenza e IDF.
    
    \item \textbf{Parametro $b \in [0, 1]$ (Penalizzazione della lunghezza documentale):}
    Regola l'incidenza della lunghezza del documento rispetto alla media della collezione:
    \begin{equation}
    B = \left(1 - b + b \cdot \frac{|D|}{\text{avgdl}}\right).
    \end{equation}
    \begin{itemize}
        \item Se $b = 1$, la normalizzazione \`e totale: un documento lungo il doppio rispetto alla media deve presentare il doppio delle occorrenze del termine per vedersi riconosciuto lo stesso score.
        \item Se $b = 0$, la normalizzazione \`e disattivata ($B = 1$ costante): la lunghezza del documento viene totalmente ignorata.
        \item Valore standard consigliato in letteratura: $b \approx 0.75$.
    \end{itemize}
\end{enumerate}

\FloatBarrier
\section{Mappa Concettuale Riassuntiva}
\label{sec:mappa-concettuale-vsm}

La Figura~\ref{fig:nlp-vsm-workflow} offre una visione d'insieme del flusso completo di elaborazione: dall'acquisizione della risorsa web grezza fino alle funzioni di scoring vettoriale e probabilistico.

\begin{figure}[H]
\centering
\resizebox{\linewidth}{!}{%
\begin{tikzpicture}[
    node distance=0.8cm and 1.2cm,
    raw/.style={rectangle, rounded corners=3pt, draw=gray!70!black, fill=gray!8,
                thick, font=\small\bfseries, minimum width=4.5cm, minimum height=0.7cm, align=center},
    stepb/.style={rectangle, rounded corners=3pt, draw=blue!70!black, fill=blue!6,
                  thick, font=\small, minimum width=4.5cm, minimum height=0.7cm, align=center},
    splitb/.style={rectangle, rounded corners=3pt, draw=teal!70!black, fill=teal!6,
                   thick, font=\footnotesize, minimum width=3.5cm, minimum height=0.7cm, align=center},
    modelb/.style={rectangle, rounded corners=3pt, draw=orange!80!black, fill=orange!8,
                   thick, font=\small\bfseries, minimum width=4.5cm, minimum height=0.7cm, align=center},
    finalb/.style={rectangle, rounded corners=3pt, draw=red!70!black, fill=red!6,
                   thick, font=\footnotesize, minimum width=4.8cm, minimum height=0.8cm, align=center},
    arr/.style={-Stealth, thick, gray!80!black}
]
    % Livello 1: Testo Grezzo
    \node[raw] (raw) {Testo Web Grezzo};

    % Livello 2: Parsing
    \node[stepb, below=of raw] (parse) {Testo Estratto\\{\scriptsize urllib + BeautifulSoup (rimozione markup)}};

    % Livello 3: Pre-processing
    \node[stepb, below=of parse] (prep) {Pulizia e Segmentazione\\{\scriptsize Tokenizzazione · Stopword Filtering · Sentence Splitting}};

    % Livello 4: Normalizzazione Morfologica (biforcazione)
    \node[splitb, below left=0.8cm and -0.8cm of prep] (stem) {\textbf{Stemming (Porter)}\\{\scriptsize Euristiche / Rapido}};
    \node[splitb, below right=0.8cm and -0.8cm of prep] (lemma) {\textbf{Lemmatizzazione (WordNet)}\\{\scriptsize Morfosintassi (POS) / Costoso}};

    % Livello 5: Vocabolario
    \node[modelb, below=2.4cm of prep] (vocab) {Vocabolario della Collezione ($V$)\\{\scriptsize Dimensione dello spazio $|V|$}};

    % Livello 6: Rappresentazioni
    \node[splitb, below left=0.8cm and -0.8cm of vocab] (bow) {\textbf{Bag of Words (BoW)}\\{\scriptsize Perde frequenza e ordine}};
    \node[splitb, below right=0.8cm and -0.8cm of vocab] (ngram) {\textbf{Estensioni N-grammi}\\{\scriptsize Word n-grams (ordine locale)}\\{\scriptsize Char n-grams (tolleranza typo)}};

    % Livello 7: VSM
    \node[modelb, below=2.4cm of vocab] (vsm) {Modello a Spazio Vettoriale (VSM)\\{\scriptsize Vettori sparsi in $\mathbb{R}^{|V|}$ (One-Hot linear combination)}};

    % Livello 8: Modelli di Scoring Finale
    \node[finalb, below left=0.8cm and -0.9cm of vsm] (tfidf) {\textbf{TF-IDF / Notazione SMART}\\{\scriptsize Log-TF + IDF ($\log(N/df)$)}\\{\scriptsize Normalizzazione Coseno ($L_2$)}};
    \node[finalb, below right=0.8cm and -0.9cm of vsm] (bm25) {\textbf{Okapi BM25}\\{\scriptsize Saturazione asintotica TF ($k_1$)}\\{\scriptsize Normalizzazione lunghezza ($b$, avgdl)}};

    % Frecce di collegamento
    \draw[arr] (raw) -- (parse);
    \draw[arr] (parse) -- (prep);
    \draw[arr] (prep.south) -- ++(0,-0.2) -| (stem.north);
    \draw[arr] (prep.south) -- ++(0,-0.2) -| (lemma.north);
    \draw[arr] (stem.south) |- (vocab.west);
    \draw[arr] (lemma.south) |- (vocab.east);
    \draw[arr] (vocab.south) -- ++(0,-0.2) -| (bow.north);
    \draw[arr] (vocab.south) -- ++(0,-0.2) -| (ngram.north);
    \draw[arr] (bow.south) |- (vsm.west);
    \draw[arr] (ngram.south) |- (vsm.east);
    \draw[arr] (vsm.south) -- ++(0,-0.2) -| (tfidf.north);
    \draw[arr] (vsm.south) -- ++(0,-0.2) -| (bm25.north);
\end{tikzpicture}%
}
\caption{Mappa concettuale integrata della pipeline di Information Retrieval: dal crawling testuale ai modelli di scoring vettoriali e probabilistici.}
\label{fig:nlp-vsm-workflow}
\end{figure}
